
All major published RLEF papers for code LLMs use binary pass/fail reward. This work proves that binary reward produces zero gradient contribution in 98.7\% of early-training GRPO groups under realistic partial-pass distributions (p\_pass = 0.1 per test case, G = 8, n = 5), and that ratio reward (k/n test cases passing) escapes this dead zone via a mathematical guarantee: whenever a GRPO group contains at least one completion with partial credit and no completion with full credit, binary reward produces zero advantage variance while ratio reward produces non-zero advantage variance (0.0000 vs. 0.0475 for the illustrative group with pass counts [0, 0, 1, 2, 0, 3, 0, 0]).

A 208-step GRPO training experiment (DeepSeek-Coder-6.7B-instruct, APPS, max\_new\_tokens = 512) produced a precision null result: both binary and ratio conditions showed reward\_mean = 0.0, clipped\_ratio = 1.0, and fraction\_partial = NaN at all 208 logged steps. The root cause is that 512-token generation is insufficient for 6.7B models to produce syntactically complete solutions for APPS competition problems, making both reward functions mathematically equivalent (both = 0 for all completions). This null result identifies the experimental prerequisite --- non-zero training-set partial solve rate --- that must be satisfied before policy-level differences between binary and ratio reward can be observed.

A complete training infrastructure (35/35 unit tests passing; 9 engineering issues documented and resolved; FractionPartialCallback diagnostic) is released to enable follow-up experiments under correctly-powered conditions. The recommended next step is re-running the policy-level experiment with max\_new\_tokens = 1,024--2,048, or using HumanEval or MBPP as training data (where DeepSeek-Coder-6.7B achieves approximately 40--65\% base pass@1). Under such conditions, the mechanistic guarantee proved here would be active, and the policy-level hypothesis becomes testable.

